% 第6章 总结与展望

\chapter{总结与展望}
\label{chap:conclusion}

\section{论文工作总结}

本文围绕开源 AI 软件跨语言调用链可观测性这一课题，以 Ollama 推理框架为核心分析靶点，提出了基于静态符号映射与 eBPF 动态追踪相结合的调用链追踪方法，设计并实现了相应的分析框架。在理论方法层面，本文提出了 Go/CGO/C++ 异构符号表五级锚点匹配方法，通过哈希等价映射、反解与 demangling 等手段解决了 CGO 边界处调用链断裂问题，覆盖 Ollama 全部 9 个 CGO 绑定函数，同时设计了动静结合的调用链缝合算法，将静态调用图（4832 节点 / 52 边）与 eBPF 追踪轨迹（1427 事件）进行融合，通过置信度计算区分 confirmed / inferred / conflicting 三类边，在包含 20 条必然调用路径的 Ground Truth 基准数据集上评估，实验结果与预期一致。在工程实现层面，本文实现了基于 eBPF 的低开销内核级追踪方案，通过将探针下沉至 llama.cpp C++ 函数层规避了 Go 1.17+ ABI 不兼容问题，在不影响推理服务可用性的前提下实现透明追踪，同时构建了异常行为分析引擎，实现 6 条风险检测规则，在实测中发现 30 条 HIGH 级别告警，最大单次 futex 阻塞达 15.2 秒，说明 AI 推理服务在并发场景下存在可用性风险。在系统架构层面，本文设计了插件化分析框架，使框架核心代码可被多个 AI 软件复用，完成了 Linux 编译环境搭建（带符号表 ollama-debug 编译）、内核验证器限制的工程化应对等基础设施工作。

\section{研究局限与不足}

本文存在以下局限性。

\textbf{CPU-only 环境，未覆盖 GPU 推理路径。} 当前追踪实验在 CPU-only 环境下进行，未覆盖 GPU 推理路径（如 CUDA kernel 调用、GPU 内存分配）。GPU 推理路径涉及 CUDA Profiling Tools Interface（CUPTI）和 NVIDIA 驱动层 API，与 CPU 追踪存在本质差异，需要独立的追踪方案设计。

\textbf{召回率受限于探针覆盖率而非算法缺陷。} 当前 confirmed 边召回率 10\%、含 inferred 边召回率 70\%，是探针覆盖率的直接体现，而非缝合算法本身的有效性限制。具体而言，52 条静态边来自 llama.cpp 源码中“所有可能执行的路径”，但实际推理请求仅走其中子集。当前仅在 8 个 llama.cpp 函数上设置 uprobe，这 8 个函数内部调用的子函数均未被探针覆盖。从结果看，10\% 召回率对应的 3 条 confirmed 边全部位于跨语言边界上，说明核心问题已被解决。

\textbf{缝合算法在极端场景下存在局限。} 对于超长调用链（>20 层）和高度动态的反射调用场景，当前基于锚点匹配的缝合算法可能失效。Go 运行时使用 \texttt{reflect} 包实现的动态方法调用无法通过静态符号表预测调用目标；C++ 虚函数的动态分派使得间接调用求解不完整。

\textbf{异常检测依赖人工定义规则。} 当前基于规则的异常检测依赖人工定义模式，对未知异常类型的泛化能力有限，缺乏对异常行为的自适应学习能力。

\section{未来工作展望}

\textbf{GPU 推理路径覆盖。} 将引入 NVIDIA GPU 追踪方案，具体包括：利用 \texttt{nvidia-smi dmon} 和 \texttt{nvidia-container-toolkit} 采集 GPU 利用率和显存使用数据；通过 CUDA Profiling Tools Interface（CUPTI）追踪 CUDA API 调用（如 \texttt{cuMemcpy}、\texttt{cuMemAlloc}），建立从 Python \texttt{torch.nn.Module.forward} 到 CUDA kernel 的完整调用链；实现容器化推理环境下的 GPU 指标采集，将容器 GPU 指标与主机侧 eBPF 追踪数据进行关联分析。

\textbf{提升探针覆盖率与缝合召回率。} 通过增加 uprobe 目标函数数量和延长采样时间，提升 confirmed 边的召回率；同时将采样窗口扩展为生产级基准数据集，捕获长尾调用路径。

\textbf{极端场景缝合优化。} 引入 Go 运行时追踪获取实际执行的反射方法名并补充到动态轨迹中；利用 RTTI 信息建立类层次结构并结合 CFG 分析求解 C++ 虚函数的可能调用目标集合；长期可引入基于深度学习的调用链预测方法。

\textbf{智能化异常检测。} 引入机器学习提升检测智能化水平，包括基于自编码器的异常检测、基于 LSTM 的时序异常检测以及构建正常行为基线，减少对人工定义规则的依赖。

\section{本章小结}

本文围绕开源 AI 推理引擎的跨语言调用链追踪问题，提出了基于 eBPF 的 HybriChain 方法，从跨语言符号映射、动态追踪采集、调用链缝合三个维度解决了混合编程架构下的调用链断裂问题。实验结果表明，HybriChain 能够在不影响推理引擎性能的前提下，完成跨语言调用链的精确追踪，并发现推理过程中的可用性脆弱点。